\documentclass[11pt]{article}
\usepackage[margin=26mm]{geometry}
\usepackage[T1]{fontenc}
\usepackage{lmodern,amsmath,amssymb,amsthm,hyperref}
\hypersetup{pdftitle={A Piecewise KL Bound from Hypergeometric Variance Envelopes},pdfauthor={Zinan Wang and Xinhao Yang}}
\newtheorem{theorem}{Theorem}
\newtheorem{proposition}[theorem]{Proposition}
\newtheorem{corollary}[theorem]{Corollary}
\newtheorem{remark}[theorem]{Remark}
\DeclareMathOperator{\Var}{Var}
\DeclareMathOperator{\Hyp}{Hyp}
\newcommand{\E}{\mathbb E}
\title{A Piecewise KL Bound from Hypergeometric Variance Envelopes}
\author{Zinan Wang\thanks{Copyright \textcopyright\ 2026 Zinan Wang. All rights reserved.}\\{\small The University of Manchester}\and Xinhao Yang\\{\small University of Southern California}}
\date{16 October 2026}
\begin{document}
\maketitle
\begin{abstract}
We integrate the pointwise minimum of two bounds on hypergeometric tilted
variance: a finite-population quadratic envelope and the classical
Poisson--binomial admissibility bound. The resulting rate has a closed
piecewise form, with a scaled Bernoulli relative entropy in its central
region and tangent-corrected binomial relative entropies in both tails.
It strictly improves the binomial KL bounds obtained by sampling and
success-count duality, and it improves the original scaled KL rate beyond
two explicit switching points. We also compare it with the optimal uniform
quadratic bound and prove that optimized factorial-moment bounds dominate
MGF-based Chernoff bounds at integer thresholds.
\end{abstract}

\section{Parameters and established variance inequalities}
Throughout Sections 1--4, assume
\begin{equation}\label{domain}
 N\ge7,\qquad 3\le m\le k\le\lfloor N/2\rfloor,\qquad N>2m.
\end{equation}
Both population parities are allowed. Let $H\sim\Hyp(N,k,m)$, and put
\[
 p=k/N,\quad q=1-p,\quad \mu=mp,\quad
 M=\frac{m(N-m)}{N-1},\quad a=\frac{N-2m}{N-2},\quad d=1-2p.
\]
All masses on $\{0,\ldots,m\}$ are positive. Write
$\Lambda(t)=\log\E e^{tH}$, $u=\Lambda'(t)$, and
$V(u)=\Lambda''(t(u))$, where $t(u)$ is the inverse tilted-mean map.
The mean map is strictly increasing from $\mathbb R$ onto $(0,m)$.

We use the two inequalities established in the accompanying manuscript
\cite{Wang}:
\begin{align}
 V(u)&\le B(u):=\frac{u(m-u)}m,\label{PB}\\
 V(u)&\le P(u):=Mpq+ad(u-\mu)-\frac{a^2}{M}(u-\mu)^2.
 \label{barrier}
\end{align}
The first follows from the classical Jacobi real-root representation as a
sum of independent Bernoulli variables; the second is its quadratic
variance-envelope refinement. In \eqref{barrier}, equality holds exactly
at $u=\mu$. In particular, $P(\mu)=Mpq>0$, and the endpoint calculations in
\cite{Wang} give $P(0)>0$ and $P(m)>0$. Concavity implies $P>0$ on $[0,m]$.
The results below are consequences and refinements of these established
inequalities; their proof does not assert a new Jacobi factorization or a
new two-level exchangeability reduction.

Define binary relative entropy by
\[
 D(z\Vert p)=z\log\frac zp+(1-z)\log\frac{1-z}{1-p},\qquad 0\le z\le1,
\]
using the continuous endpoint values. All logarithms are natural.

\section{Integration in the tilted mean}
\begin{theorem}[Variance-envelope integration]\label{integration}
Let $X$ have finite support with minimum $l$, maximum $h$, and positive
mass at both endpoints, with $l<h$. Let $\mu=\E X$ and
$I(v)=\sup_t\{tv-\log\E e^{tX}\}$.
Suppose its tilted variance $V(u)$ is bounded by a continuous function
$Q(u)>0$ on $(l,h)$. Set
\[
 J_Q(v)=
 \begin{cases}
 \displaystyle\int_\mu^v\frac{v-u}{Q(u)}\,du,&\mu\le v<h,\\[5pt]
 \displaystyle\int_v^\mu\frac{u-v}{Q(u)}\,du,&l<v\le\mu.
 \end{cases}
\]
Then $I(v)\ge J_Q(v)$. Consequently
\[
 \Pr(X\ge v)\le e^{-J_Q(v)}\quad(v\ge\mu),\qquad
 \Pr(X\le v)\le e^{-J_Q(v)}\quad(v\le\mu).
\]
The same statements hold at $l,h$ whenever the displayed integrals have
finite one-sided limits. If $V<Q$ on an interval of positive length between
$\mu$ and $v$, then $I(v)>J_Q(v)$.
\end{theorem}
\begin{proof}
The tilted-mean map is strictly increasing, with derivative $V>0$.
For $v\in(l,h)$, differentiation of its Legendre transform gives
$I'(v)=t(v)$, $I''(v)=1/V(v)$, and $I(\mu)=I'(\mu)=0$.
Integrating twice gives exactly the displayed integrals with $V$ in place
of $Q$. Their kernels are nonnegative on the respective integration
intervals, proving both the comparison and its strict version.
Markov's inequality, optimized over $t\ge0$ for the upper tail and
$t\le0$ for the lower tail, proves the tail assertions.
At the endpoints, the finite-support exponential sum gives
$I(l)=-\log\Pr(X=l)$ and $I(h)=-\log\Pr(X=h)$.
Taking limits in the integral comparisons completes the proof.
\end{proof}

The rate obtained from $P$ is
\begin{equation}\label{IP}
 I_P(v)=\frac M{a^2}D\left(p+\frac aM(v-\mu)\,\middle\Vert\,p\right).
\end{equation}
Indeed, $I_P''=1/P$ and $I_P(\mu)=I_P'(\mu)=0$.
Its KL argument remains strictly inside $(0,1)$ for $0\le v\le m$, since
\[
 1-\frac{am}{M}
 =\frac{N(m-1)}{(N-2)(N-m)}>0.
\]
Thus an endpoint of the hypergeometric support never requires an invalid
KL argument in \eqref{IP}. The corresponding rate for $B$ is
\begin{equation}\label{IB}
 I_B(v)=mD(v/m\Vert p).
\end{equation}

For completeness, direct optimization of the established Laplace bound
\[
 \log\E e^{t(H-\mu)}\le
 \frac M{a^2}\{\log(q+pe^{at})-pat\}
\]
at threshold $v$ gives
$t_*=a^{-1}\log[\theta q/(p(1-\theta))]$,
where $\theta=p+a(v-\mu)/M$, and exponent $I_P(v)$.
This is the ordinary Bernoulli Chernoff transform. For an abstract
Bernoulli envelope, $\theta=0,1$ are interpreted by $t\to\mp\infty$;
values outside $[0,1]$ are not inserted into $D$.

\section{A closed piecewise rate}
Put
\[
 c=\frac1m-\frac{a^2}{M},\qquad
 h=(1-a)d,\qquad b=(m-M)pq,
\]
and define
\begin{equation}\label{switch}
 r_\pm=\mu+\frac{h\pm\sqrt{h^2+4cb}}{2c}.
\end{equation}
\begin{theorem}[Piecewise KL rate]\label{piecewise}
Under \eqref{domain}, $0<r_-<\mu<r_+<m$. The rate obtained by integrating
$1/\min(P,B)$ is
\begin{equation}\label{J}
 J(v)=
 \begin{cases}
 I_B(v)+I_P(r_-)-I_B(r_-)
 +(v-r_-)[I_P'(r_-)-I_B'(r_-)],&0\le v<r_-,\\
 I_P(v),&r_-\le v\le r_+,\\
 I_B(v)+I_P(r_+)-I_B(r_+)
 +(v-r_+)[I_P'(r_+)-I_B'(r_+)],&r_+<v\le m.
 \end{cases}
\end{equation}
It gives valid upper and lower tail bounds as in Theorem~\ref{integration}.
For every $v\ne\mu$ in $[0,m]$,
\begin{equation}\label{strict}
 I(v)>J(v)>I_B(v),\qquad J(v)\ge I_P(v),
\end{equation}
and $J(v)>I_P(v)$ holds exactly outside $[r_-,r_+]$.
\end{theorem}
\begin{proof}
Writing $y=u-\mu$ gives $P(u)-B(u)=cy^2-hy-b$.
Here $b>0$ and
\[
 c=\frac{(m-1)\{3N^2-4N-4m(N-1)\}}
 {m(N-m)(N-2)^2}>0,
\]
because $m<N/2$ makes the bracket exceed $N(N-2)>0$.
The difference is negative at $\mu$ and positive at $0,m$; hence its
two roots are precisely \eqref{switch}, with the asserted locations.
It follows that $\min(P,B)=P$ between the roots and $B$ outside.
Integrating with zero value and derivative at $\mu$ yields \eqref{J};
the affine corrections are forced by continuity of the rate and its
first derivative at each root. In fact $J$ is twice continuously
differentiable at the roots because $P=B$ there.

Comparison of the nonnegative integration kernels gives $J\ge I_P,I_B$.
On every noncentral path the interval near $\mu$ has $P<B$, proving
$J>I_B$. A path leaving $[r_-,r_+]$ also crosses an interval where $B<P$,
proving the characterization of $J>I_P$. Finally, $V<P$ except at $\mu$,
so the central interval gives $I>J$ on every noncentral path.
At $0,m$, $1/B$ has simple poles, canceled by the endpoint integration
kernels; the limits are finite and the same strict arguments apply.
\end{proof}

\begin{corollary}[Integer and two-sided events]
For real $x>0$, set
$z_+=\lceil\mu+x\rceil$ and $z_-=\lfloor\mu-x\rfloor$.
Use $e^{-J(z_+)}$ for the upper event when $z_+\le m$, and zero otherwise;
use $e^{-J(z_-)}$ for the lower event when $z_-\ge0$, and zero otherwise.
Their sum, capped at one, bounds $\Pr(|H-\mu|\ge x)$.
\end{corollary}

\section{Comparison boundaries}
\subsection{Four binomial KL bounds}
Sampling duality and complementation give the following four classical
binomial Chernoff rates at a fixed threshold $v$:
\begin{align*}
 R_1(v)&=mD(v/m\Vert k/N),&
 R_2(v)&=kD(v/k\Vert m/N),\\
 R_3(v)&=(N-m)D((k-v)/(N-m)\Vert k/N),&
 R_4(v)&=(N-k)D((m-v)/(N-k)\Vert m/N).
\end{align*}
All have value and derivative zero at $\mu$. Their reciprocal second
derivatives are, respectively,
\[
 B_1=B,\quad B_2=\frac{v(k-v)}k,\quad
 B_3=\frac{(k-v)(N-m-k+v)}{N-m},\quad
 B_4=\frac{(m-v)(N-k-m+v)}{N-k}.
\]
With $L=k-m\ge0$ and $R=N-m-k\ge0$,
\begin{align*}
 B_2-B&=\frac{(k-m)v^2}{mk}\ge0,\\
 B_4-B&=\frac{(N-k-m)(m-v)^2}{m(N-k)}\ge0,\\
 B_3-B&=\frac{Lv^2+R(m-v)^2+mLR}{m(N-m)}\ge0.
\end{align*}
Therefore $R_1=I_B\ge R_j$ for $j=2,3,4$ on $[0,m]$, and
Theorem~\ref{piecewise} strictly improves their maximum at every
noncentral threshold. Duality itself is classical and is explicitly used
by Mannalath, Zapatero and Curty~\cite{MZC}; the improvement here comes
from integrating the smaller mean-dependent variance bound.

\subsection{The optimal uniform quadratic coefficient}
Define
\[
 \kappa_N=\begin{cases}1,&N\text{ even},\\
 [N\operatorname{artanh}(1/N)]^{-1},&N\text{ odd},\end{cases}
 \qquad V_* = M\kappa_N/4.
\]
This is the coefficient for $H-\mu$ obtained from the exact uniform
exchangeability constant in \cite{Wang}, in the convention
$\log\E e^{t(H-\mu)}\le V_*t^2/2$.
The optimal Bernoulli proxy is
\[
 \beta(p)=\frac{1-2p}{2\log((1-p)/p)},\qquad \beta(1/2)=1/4.
\]
Under \eqref{domain}, $M\beta(p)\le V_*$. Consequently the whole scaled
Laplace envelope is bounded above by $M\beta(p)t^2/2$, and conjugation gives
\begin{equation}\label{quadratic}
 J(v)\ge I_P(v)\ge \frac{(v-\mu)^2}{2M\beta(p)}
 \ge\frac{(v-\mu)^2}{2V_*}.
\end{equation}
Thus the scaled KL and piecewise bounds are never weaker than this
particular uniform quadratic bound. They need not dominate the numerically
optimized quadratic proxy for each fixed hypergeometric law.

For $p<1/2$, equality in the middle inequality at a nonzero displacement
requires $\theta=1-p$, hence $v-\mu=M(1-2p)/a>0$; this follows from the
unique nonzero equality in the Bernoulli proxy inequality.
For odd $N$, equality with the final uniform rate additionally requires
$k=(N-1)/2$. For even $N$, equality $M\beta(p)=V_*$ requires $p=1/2$,
but the balanced Bernoulli quadratic inequality is strict at nonzero tilt.
The inequality $J\ge I_P$ is an equality precisely in the central region.

Writing $x=v-\mu$, the local expansion is
\[
 I_P(\mu+x)=\frac{x^2}{2Mpq}
 -\frac{a(1-2p)x^3}{6M^2p^2q^2}
 +\frac{a^2(p^3+q^3)x^4}{12M^3p^3q^3}+O(x^5).
\]
This shows the directional asymmetry and the variance-scale improvement
near the mean. Since $M<m$, $I_P>I_B$ locally, but the unspliced $I_P$
can be smaller than $I_B$ in the tails. Theorem~\ref{piecewise} removes
this loss relative to the four binomial bounds.

\subsection{Exact-variance Bennett and Bernstein bounds}
Put $v_0=Mpq$ and $h_0(z)=(1+z)\log(1+z)-z$.
The Bennett exponent in Greene--Wellner~\cite[Theorem 3]{GW} is
$v_0h_0(x/v_0)$; their Bernstein corollary replaces it by
$x^2/[2(v_0+x/3)]$. For either tail, the scaled rate is strictly larger
than the Bennett exponent at every nonzero supported deviation.
Indeed, for $s>0$,
\[
 P(\mu+s)=v_0+ad s-a^2s^2/M<v_0+s,\qquad
 P(\mu-s)=v_0-ad s-a^2s^2/M<v_0+s,
\]
because $0\le ad<1$. The second derivative of $v_0h_0(x/v_0)$ is
$1/(v_0+x)$, while that of $I_P(\mu\pm x)$ is $1/P(\mu\pm x)$.
Both rates and their first derivatives vanish at zero. Integrating the
strict comparison twice, and then using $J\ge I_P$, proves the assertion.
The standard inequality $h_0(z)\ge z^2/[2(1+z/3)]$ gives the Bernstein
comparison as well. This does not compare $J$ with Greene--Wellner's
other bounds containing a prefactor and a quartic exponent.

\subsection{Optimized factorial moments}
\begin{proposition}[A comparison with exact Chernoff]\label{factorial}
Let $X$ be nonnegative and integer-valued with bounded support, and let
$d\ge1$ be an integer. Define
\[
 F_d=\min_{0\le j\le d}\frac{\E\binom Xj}{\binom dj}.
\]
Then
\[
 \Pr(X\ge d)\le F_d\le
 \inf_{t\ge0}\E e^{t(X-d)}.
\]
For the hypergeometric law,
\[
 \frac{\E\binom Hj}{\binom dj}
 =\frac{(k)_{\underline j}(m)_{\underline j}}
 {(N)_{\underline j}(d)_{\underline j}},\qquad 0\le j\le d.
\]
In particular, at integer upper thresholds $d>\mu$ in \eqref{domain},
\[
 \Pr(H\ge d)\le F_d\le e^{-I(d)}<e^{-J(d)}.
\]
The lower-tail version follows by applying the result to $m-H$.
\end{proposition}
\begin{proof}
For $j\le d$, $\binom Xj\ge\binom dj$ on $\{X\ge d\}$, which proves
the first inequality. For $z=e^t-1\ge0$, the binomial theorem gives
\[
 \E(1+z)^X=\sum_{j\ge0}\E\binom Xj z^j
 \ge F_d\sum_{j=0}^d\binom dj z^j=F_d(1+z)^d.
\]
Divide by $e^{td}$ and take the infimum. Hypergeometric factorial
moments give the displayed ratio.
\end{proof}

Optimized hypergeometric factorial-moment bounds are already studied
in \cite{MZC}; neither that method nor its optimization is new here.
Earlier comparisons of moment bounds with Chernoff bounds include
Philips and Nelson~\cite{PN}; Proposition~\ref{factorial} provides the
finite integer-valued binomial-basis argument used for the present comparison.
For $\mu<d\le\min(k,m)$, a valid minimizing index is
\[
 j_* =\left\lceil\frac{Nd-km}{N+d-k-m}\right\rceil\in\{1,\ldots,d\}.
\]
Indeed the ratio of consecutive candidates is at most one exactly when
$j\le(Nd-km)/(N+d-k-m)$. When this fraction is integral, the two adjacent
indices tie if both are admissible. At $d=\min(k,m)$ the valid endpoint
index is $j=d$; one must not use the undefined denominator
$(d)_{\underline{d+1}}$. These facts set a precise comparison boundary:
the piecewise KL expression is an explicit refinement of the stated KL
and uniform quadratic envelopes, not a universal improvement over
optimized factorial moments or all existing finite-population bounds.

\begin{thebibliography}{9}
\bibitem{Wang} Z. Wang, \emph{Exact Hoeffding Constants for Finite Exchangeable
Sequences}, 2026, main repository source at commit
\texttt{ebf534362f596175a24fc33e9eebc1cd4145349f},
the uniform Laplace comparison and quadratic variance-envelope lemma in
\texttt{source/odd\_uniform.tex}.
\bibitem{Hoeffding} W. Hoeffding, Probability inequalities for sums of bounded
random variables, \emph{J. Amer. Statist. Assoc.} 58 (1963), 13--30.
\bibitem{Chvatal} V. Chv\'atal, The tail of the hypergeometric distribution,
\emph{Discrete Math.} 25 (1979), 285--287.
\bibitem{MZC} V. Mannalath, V. Zapatero and M. Curty,
\emph{Elementary tail bounds on the hypergeometric distribution},
arXiv:2510.19726v2 (2025), Propositions 1--2;
\url{https://arxiv.org/abs/2510.19726v2}; published in
\emph{Communications in Statistics---Theory and Methods} 55 (2026),
6938--6943, \url{https://doi.org/10.1080/03610926.2026.2676012}.
\bibitem{GW} E. Greene and J. A. Wellner, Exponential bounds for the
hypergeometric distribution, \emph{Bernoulli} 23 (2017), 1911--1950.
\bibitem{BM} R. Bardenet and O.-A. Maillard, Concentration inequalities for
sampling without replacement, \emph{Bernoulli} 21 (2015), 1361--1385.
\bibitem{PN} T. K. Philips and R. Nelson, The moment bound is tighter than
Chernoff's bound for positive tail probabilities,
\emph{The American Statistician} 49 (1995), 175--178.
\end{thebibliography}
\end{document}
